\chapter{Current Memory Reduction Methods}
\label{ch:current-memory-reduction-methods}

Recent all-atom generative models such as RFDiffusion 3~\cite{butcher_novo_2025} (RFD3), 
which is inspired by the AlphaFold 3 architecture~\cite{abramson_accurate_2024} (AF3),
can in principle generate multimeric structures such as the symmetric assembly needed as scaffolds
for nanoparticle vaccines. This is achieved by jointly modelling all chains with
an atomistic level of precision and designing proteins with predefined point-group 
symmetries. However, the underlying architecture maintains pairwise representations 
whose memory cost scales quadratically with the number of tokens~$I$ (and atoms~$L$). 
For protein assemblies, $I$ and $L$ grow linearly with the number of chains
modelled, causing the $\mathcal{O}$($I^{2}$) and $\mathcal{O}$($L^{2}$) pairwise 
feature tensors and related quadratic intermediates to exceed the memory capacity 
of a single GPU when trying to model assemblies of the size relevant for nanoparticle vaccines. 
This practical bottleneck therefore heavily limits the size of what can be
designed.

% This single-device ceiling, however, contrasts with broader trends in 
% computational infrastructure: while single-device memory capacity has grown only 
% incrementally, access to multi-GPU clusters is now routine for both academic and 
% industrial research groups. Accordingly, partitioning large transformer workloads 
% across such clusters has become standard practice in adjacent fields such as large 
% language modelling, both for training~\cite{li_sequence_2022} and 
% inference~\cite{pope_efficiently_2022}.

\section{Reducing the cost of attention}
\label{sec:reducing-attention-cost}

In the literature (especially in the one related to large language models), several strategies have 
been proposed to reduce the memory cost of attention. Here we briefly cover 
IO-efficient attention on a single device, while devoting most of this sec-
tion to distributed parallelism for structure models, because our developed solution 
\cref{ch:methods} sits in this realm. 

\subsection{IO-efficient attention}
\label{subsec:io-efficient-attention}
FlashAttention \cite{dao_flashattention_2022, dao_flashattention-2_2023}
computes exact attention with memory linear in sequence length 
($\mathcal{O}(I)$ or $\mathcal{O}(L)$ in our notation) by tiling queries, keys, 
and values into SRAM-resident blocks while maintaining running softmax statistics 
\citep{milakov_online_2018}, building on the earlier observation that attention 
admits a linear-memory implementation \citep{rabe_self-attention_2022}. 
For architectures without pair representations, this suffices and motivates its 
use in protein language models like ESM3 \citep{hayes_simulating_2025}, which 
represent sequence, structure, and function as discrete tokens and condense 
pairwise geometry into a single SE(3)-invariant geometric-attention block at 
the input. Flash‑IPA \citep{liu_flash_2025} and FlashBias \citep{wu_flashbias_2025} 
extend this idea to pair‑biased attention, such as the one present in AlphaFold‑3 
Pairformer, by re-expressing geometric and pairwise bias terms via additional low‑rank 
features concatenated into the query and key projections. For complex learned pair
biases, these methods generally require training auxiliary parameters and are not 
a weight‑preserving drop‑in for an arbitrary pretrained model. 
% These approaches are 
% orthogonal to Design-CP: they reduce the per-block memory footprint on a single device, 
% while we shard the persistent quadratic activations across devices, and the two could in principle be composed.

\subsection{Distributed parallelism for structure models.} 
Among the standard parallelism axes (data, tensor, pipeline, expert, activation), 
\emph{context parallelism} uniquely shards activations along the sequence dimension 
of every layer, which makes it a natural fit for the $\mathcal{O}(I^2)$ pair tensor 
of AlphaFold-class models. FastFold \citep{cheng_fastfold_2023} introduced Dynamic 
Axial Parallelism (DAP) for AlphaFold~2, replicating parameters on every device and 
sharding activations along a single sequence axis at a time; 
ScaleFold \citep{zhu_scalefold_2024} adopted DAP and scaled to 2048 H100s for training. 
As the Evoformer interleaves row- and column-wise attention over the MSA track, 
DAP must insert an all-to-all communication step whenever the active axis flips, 
incurring six all-to-all redistributions per Evoformer block at inference, together
with one \textsc{AllGather} in the outer-product-mean and two in the triangular
updates~\citep{cheng_fastfold_2023}. These collectives contribute non-trivially 
to inference latency at scale and increase transient memory relative to the 
steady-state ($\mathcal{O}(I^2/P)$, where $P$ is the total number of devices) sharded pair representation.

Fold-CP \citep{lin_fold-cp_2026} generalises axis sharding to a full two-dimensional 
context-parallel strategy for AF3-class models (implementing it for 
Boltz-2~\citep{passaro_boltz-2_2025}), extending Ring Attention~\citep{liu_ring_2023} 
into a Cannon-style 2D ring tailored to dense triangular updates, while window-batched
atom attention is handled by a complementary shardwise kernel that keeps each 
window's attention local to its rank. The pair tensor is tiled over a 
$\sqrt{P}\times\sqrt{P}$ device grid, so that rank $(r,c)$ holds the block indexed 
by residue ranges $[r\cdot I/\sqrt{P}, (r{+}1)\cdot I/\sqrt{P}]\times [c\cdot I/\sqrt{P}, (c{+}1)\cdot I/\sqrt{P}]$. 
Triangle attention, triangle multiplication, attention-with-pair-bias, pair-weighted-averaging, 
and outer-product-mean are each reformulated as ring algorithms in which $\mathbf{K}$ 
and $\mathbf{V}$ shards (and, where required, triangular biases and masks) circulate 
between neighbouring devices while a numerically-stable tiled softmax merges partial 
outputs without ever materialising the full $I\times I$ attention on any rank. 
The resulting steady-state pair memory per device is ($\mathcal{O}(I^2/P)$), matching 
1D axis-sharded DAP. However, under 1D sharding, triangular updates typically require 
collectives over all ($P$) ranks (e.g., to obtain the necessary key/value or bias 
shards along the active axis), which can inflate the transient working-set memory 
during attention/multiplication beyond the steady shard. In Fold‑CP’s 2D tiling, 
collectives are restricted to a single row or column subgroup of size ($\sqrt{P}$). 
This reduces communication volume and confines the transient working-set memory growth
in triangular updates (e.g., gathered/circulated K/V shards, triangle biases, masks,
and other pair-like intermediates) to ($\sqrt{P}$)-sized groups, rather than requiring 
global ($P$)-way collectives.
